% !TEX root = ../main.tex
% !TEX program = lualatex
\documentclass[../main.tex]{subfiles}
\IfSubfilesClassLoaded{\externaldocument{\subfix{../build/main}}}{}
% =============================================================================
% File: 27H_High_Consequence_Events_Management.tex
% =============================================================================
\begin{document}
\IfSubfilesClassLoaded{\chapter{EDO Study Guide}}{}
\section{High Consequence Event Management}

\subsection{Summary}
\begin{description}
  \item[\textbf{Risk-aware culture.}] A deliberate culture that promotes risk-aware behaviors and confronts human element weaknesses reduces the likelihood of high consequence events~\autocite{edo-5-1-10-hce-management-2025}.
  \item[\textbf{Human element focus.}] Technical strength alone is insufficient; human element weaknesses often enable failures and must be addressed explicitly~\autocite{edo-5-1-10-hce-management-2025}.
  \item[\textbf{HCEPF alignment.}] The \ac{hcepf} links risk-aware behaviors to Get Real, Get Better and learning teams to sustain safe, effective performance~\autocite{edo-5-1-10-hce-management-2025}.
\end{description}

\subsection{Why High Consequence Event Prevention}
\begin{description}
  \item[\textbf{Culture first.}] Technically strong organizations still fail when cultural flaws go unaddressed; a strong culture plus technical foundation preempts catastrophic mistakes~\autocite{edo-5-1-10-hce-management-2025}.
  \item[\textbf{Fukushima lesson.}] The Fukushima disaster illustrates how human and cultural factors compound technical risk and drive catastrophic outcomes~\autocite{edo-5-1-10-hce-management-2025}.
\end{description}

\subsection{Risk-Aware Culture}
\begin{description}
  \item[\textbf{Definition.}] A risk-aware culture uses daily risk-aware behaviors, expects accountability for human element weaknesses, and empowers teams to identify and mitigate mission risks~\autocite{edo-5-1-10-hce-management-2025}.
  \item[\textbf{Deliberate design.}] Organizational habits must be intentionally designed rather than left to chance~\autocite{edo-5-1-10-hce-management-2025}.
\end{description}

\subsection{Heinrich's Triangle}
\begin{description}
  \item[\textbf{Human error link.}] Major accidents often share the same causal human element weaknesses seen in minor accidents and near-misses~\autocite{edo-5-1-10-hce-management-2025}.
  \item[\textbf{Prevention strategy.}] Track and correct small failures to reduce the likelihood of apex events~\autocite{edo-5-1-10-hce-management-2025}.
\end{description}

\subsection{Performance Focus}
\begin{description}
  \item[\textbf{Three domains.}] Effective risk management requires attention to technology, systems/processes, and people behaviors; overreliance on any single domain creates gaps~\autocite{edo-5-1-10-hce-management-2025}.
  \item[\textbf{Barriers.}] Programmatic system barriers can fail when risks align and decision errors defeat safeguards~\autocite{edo-5-1-10-hce-management-2025}.
\end{description}

\subsection{High Consequence Events Prevention Framework}
\begin{description}
  \item[\textbf{Framework scope.}] \ac{hcepf} is a structured cultural initiative with 24 risk-aware behaviors, each paired to a human element weakness to track learning and performance~\autocite{edo-5-1-10-hce-management-2025}.
  \item[\textbf{Get Real, Get Better linkage.}] \ac{hcepf} behaviors are the mechanism for getting real, getting better, and building learning teams within \ac{ssp}~\autocite{edo-5-1-10-hce-management-2025}.
  \item[\textbf{Tools.}] The \ac{hcepf} tools reference guide provides a menu of practices to strengthen behaviors at the individual, team, and organizational levels~\autocite{edo-5-1-10-hce-management-2025}.
\end{description}

\subsection{Balancing Risk Aversion}
\begin{description}
  \item[\textbf{Competing pressures.}] Production, cost, and schedule pressures compete with effectiveness, security, and safety; leaders must manage the tension and avoid drifting toward unsafe extremes~\autocite{edo-5-1-10-hce-management-2025}.
\end{description}

\subsection{Key Finding}
\begin{description}
  \item[\textbf{Technical + human strengths.}] Organizations with cross-functional technical skill and strong human element behaviors prevent high consequence events and respond with resilience when surprises occur~\autocite{edo-5-1-10-hce-management-2025}.
  \item[\textbf{Root cause discipline.}] Technical errors must be traced to underlying human element causes to avoid recurrence~\autocite{edo-5-1-10-hce-management-2025}.
\end{description}

\subsection{Case Study Prompts}
\begin{itemize}
  \item Evaluate human element weaknesses in real-world incidents (e.g., OceanGate Titan, Kobe Bryant helicopter crash, Apollo 13 oxygen tank specification error, Grenfell Tower) to identify behaviors to strengthen~\autocite{edo-5-1-10-hce-management-2025}.
\end{itemize}

\subsection{Quick Review}
\begin{itemize}
  \item A risk-aware culture is deliberately designed, reinforced daily, and essential to high consequence event prevention~\autocite{edo-5-1-10-hce-management-2025}.
  \item Human element weaknesses enable technical failures to cascade into disasters~\autocite{edo-5-1-10-hce-management-2025}.
  \item \ac{hcepf} behaviors link to Get Real, Get Better and provide tools for continuous improvement~\autocite{edo-5-1-10-hce-management-2025}.
\end{itemize}

\IfSubfilesClassLoaded{
  \RenewDocumentCommand{\entryname}{}{\textbf{\color{Modern} Acronym}}
  \RenewDocumentCommand{\descriptionname}{}{\textbf{\color{Modern} Definition}}
  \printnoidxglossary[
    type=\acronymtype,
    title=Acronyms,
    style=long-booktabs
  ]}{}
\end{document}
